\section{absolute convergence}

For series with positive terms, we have many convergence criteria that, in
general, do not apply to series with arbitrary signs or to series of complex
numbers. However, the convergence of the latter can sometimes be easily reduced
to the convergence of series with positive terms by considering the modulus of
each term.

\begin{definition}[absolute convergence]
\mymark{***}
We say that a series (with real or complex terms) $\sum a_n$ is \emph{absolutely
convergent}%
\mymargin{absolutely convergent}%
\index{absolutely!convergent} if the series $\sum \abs{a_n}$ is convergent.
\end{definition}

\begin{theorem}[absolute convergence]\label{th:convergenza_assoluta}
\mymark{***}%
If a series $\sum a_n$ (real or complex) is absolutely convergent, then it is
convergent and the following holds:
\[
  \abs{\sum_{k=0}^{+\infty} a_k} \le \sum_{k=0}^{+\infty} \abs{a_k}.
\]
\end{theorem}
%
\begin{proof}
\mymark{*}
Initially, suppose that the $a_n$ are real numbers. Define $a_n^+ =
\max\ENCLOSE{0, a_n}$ and $a_n^- = -\min \ENCLOSE{0, a_n}$. That is, if $a_n\ge
0$, then $a_n^+ = a_n$ and $a_n^-=0$; if $a_n\le 0$, then $a_n^+ =0$ and $a_n^-
= -a_n$. Thus, $a_n^+\ge 0$, $a_n^-\ge 0$,
\[
   a_n = a_n^+  - a_n^-
   \qquad\text{and}\qquad
   \abs{a_n} = a_n^+ + a_n^-.
\]
Then, if $\sum \abs{a_n}$ converges, by comparison, both $\sum a_n^+$ and $\sum
a_n^-$ converge. Therefore, by the theorem on the sum of limits, $\sum a_n =
\sum a_n^+ - \sum a_n^-$, and hence $\sum a_n$ also converges.

If we have a sequence of complex numbers $a_n = x_n + i y_n$ and if $\sum
\abs{a_n}$ converges, then by comparison, both $\sum \abs{x_n}$ and
$\sum\abs{y_n}$ converge (note that $\abs{x} \le \abs{x+iy}$ and $\abs{y}\le
\abs{x+iy}$). Thus, $\sum x_n$ and $\sum y_n$ converge by what has already been
demonstrated for series with real terms. But then $\sum i y_n$ and $\sum a_n =
\sum (x + iy_n)$ also converge.

Now, let us set
\[
  S_n  = \sum_{k=0}^n a_k.
\]
By the subadditivity of the modulus, we know that for finite sums, we have
\[
 \abs{S_n} =\abs{\sum_{k=0}^n a_k}
 \le \sum_{k=0}^n \abs{a_k} \le \sum_{k=0}^{+\infty} \abs{a_k}.
\]
And by the continuity of the modulus, setting $S= \lim S_n$, we have
\[
  \abs{\sum_{k=0}^{+\infty} a_k}
  = \abs{S}
  = \lim_{n\to +\infty} \abs{S_n}
  \le \sum_{k=0}^{+\infty} \abs{a_k}.
\]
\end{proof}

\begin{theorem}[exchange of series]
\label{th:scambio_somma}
Let $a_{k,j}\in \RR$ or $a_{k,j}\in \CC$ be a two-index sequence $k\in \NN$,
$j\in \NN$. If
\begin{equation}\label{eq:499655}
  \sum_{k=0}^{+\infty}\sum_{j=0}^{+\infty}\abs{a_{k,j}}<+\infty  
  \qquad \text{or} \qquad
  \sum_{j=0}^{+\infty}\sum_{k=0}^{+\infty}\abs{a_{k,j}}<+\infty  
\end{equation}
then
\begin{equation}\label{eq:scambio_somma}
  \sum_{k=0}^{+\infty} \sum_{j=0}^{+\infty} a_{k,j}
  = \sum_{j=0}^{+\infty} \sum_{k=0}^{+\infty} a_{k,j}.
\end{equation}
\end{theorem}
%
\begin{proof}
Since the statement is symmetric in $k$ and $j$, we can assume that the first of
the two alternatives in hypothesis~\eqref{eq:499655} is satisfied: $\sum_k
\sum_j \abs{a_{k,j}}<+\infty$.

It is understood, of course, that for each $k$, the series $\sum_j
\abs{a_{k,j}}$ is convergent, hence $\sum_j a_{k,j}$ is absolutely convergent.
% and, by theorem~\ref{th:convergenza_assoluta}, 
% also $\sum_k \abs{\sum_j a_{k,j}} 
% \le \sum_k \sum_j \abs{a_{k,j}}$
% is absolutely convergent.
Obviously, 
$\sum_k \abs{a_{k,j}} \le \sum_k \sum_j \abs{a_{k,j}}$
and thus for each $j$, the series 
$\sum_k a_{k,j}$ is also absolutely convergent.
Setting
\[
S = \sum_{k=0}^{+\infty}\sum_{j=0}^{+\infty} a_{k,j},
\qquad
S_n = \sum_{j=0}^{n-1}\sum_{k=0}^{+\infty} a_{k,j}  
\]
we need to show that $S_n \to S$
as $n\to +\infty$. 
We apply the definition of limit:
let $\eps>0$ be fixed.
By the tail theorem~\ref{th:coda}, there exists $K$ such that 
\[
  \sum_{k=K}^{+\infty} \sum_{j=0}^{+\infty} \abs{a_{k,j}}<\frac \eps 2
\]
and for each $k<K$, there exists $J_k$ such that 
\[
  \sum_{j=J_k}^{+\infty} \abs{a_{k,j}} < \frac{\eps}{2K}.  
\]
Thus, setting $J=\max\ENCLOSE{J_0,J_1, \dots, J_{K-1}}$, if $n>J$, we have:
\begin{align*}
  \abs{S-S_n}
  &= \abs{\sum_{k=0}^{+\infty}\sum_{j=0}^{+\infty}a_{k,j} 
  - \sum_{j=0}^{n-1}\sum_{k=0}^{+\infty}a_{k,j}}
  = 
  \abs{\sum_{k=0}^{+\infty}\sum_{j=0}^{+\infty}a_{k,j} 
  - \sum_{k=0}^{+\infty}\sum_{j=0}^{n-1}a_{k,j}}\\
  &=
  \abs{\sum_{k=0}^{+\infty}\sum_{j=n}^{+\infty}a_{k,j} }
  = \abs{\sum_{k=0}^{K-1}\sum_{j=n}^{+\infty} a_{k,j}
   + \sum_{k=K}^{+\infty}\sum_{j=n}^{+\infty} a_{k,j}}\\
   &\le \sum_{k=0}^{K-1}\sum_{j=n}^{+\infty} \abs{a_{k,j}}
    + \sum_{k=K}^{+\infty}\sum_{j=0}^{+\infty} \abs{a_{k,j}}
   \le \sum_{k=0}^{K-1}\sum_{j=J_k}^{+\infty} \abs{a_{k,j}}
   + \frac \eps 2
  \le \frac \eps 2 + \frac \eps 2 = \eps
\end{align*}
as we wanted to demonstrate.
\end{proof}

\begin{theorem}[unconditional convergence]%
\label{th:convergenza_incondizionata}%
\mymark{*}%
\index{convergence!unconditional}%
If $\sum a_k$ is an absolutely convergent series and $\sigma\colon \NN \to \NN$
is any bijective function (permutation of natural numbers), then
\[
  \sum_{k=0}^{+\infty} a_k
  = \sum_{j=0}^{+\infty} a_{\sigma(j)}.
\]
\end{theorem}
\begin{proof}
Define the two-index sequence $b_{k,j}$ as follows:
\[
  b_{k,j} = \begin{cases}
    a_k & \text{if $k=\sigma(j)$},\\
    0 & \text{otherwise}.
  \end{cases}  
\] 
Clearly, $\sum_j \abs{b_{k,j}} = \abs{a_k}$, and thus $\sum_k \sum_j
\abs{b_{k,j}} = \sum_k \abs{a_k} < +\infty$. Therefore, by
theorem~\ref{th:scambio_somma}
\[
 \sum_{k=0}^{+\infty} a_k 
 = \sum_{k=0}^{+\infty} \sum_{j=0}^{+\infty}b_{k,j}
 = \sum_{j=0}^{+\infty} \sum_{k=0}^{+\infty}b_{k,j}
 = \sum_{j=0}^{+\infty} a_{\sigma(j)}.  
\]
\end{proof}

\begin{theorem}[Cauchy sums]
  \label{th:somma_Cauchy}%
Let $a_{k,j}\in \RR$ or $a_{k,j}\in \CC$ be a two-index sequence $k\in \NN$,
$j\in\NN$. If
\[
  \sum_{k=0}^{+\infty} \sum_{j=0}^{+\infty} \abs{a_{k,j}} < +\infty
  \qquad\text{or}\qquad 
  \sum_{n=0}^{+\infty} \sum_{k=0}^n \abs{a_{k,n-k}} < +\infty
\]  
then 
\begin{equation}\label{eq:somma_Cauchy}
  \sum_{k=0}^{+\infty} \sum_{j=0}^{+\infty} a_{k,j}
   = \sum_{n=0}^{+\infty} \sum_{k=0}^{n} a_{k,n-k}.
\end{equation}
\end{theorem}
%
\begin{proof}
Set 
\[
  b_{k,n} = \begin{cases}
    a_{k,n-k} & \text{if $k\le n$}\\
    0 & \text{if $k>n$}.
  \end{cases}  
\]
Then observe that 
\begin{align*}
  \sum_{k=0}^{+\infty} \sum_{n=0}^{+\infty} b_{k,n}
  &= \sum_{k=0}^{+\infty} \sum_{n=k}^{+\infty} a_{k,n-k} 
  = \sum_{k=0}^{+\infty} \sum_{j=0}^{+\infty} a_{k,j}
  \\
  \sum_{n=0}^{+\infty}\sum_{k=0}^{+\infty} b_{k,n} 
  &= \sum_{n=0}^{+\infty}\sum_{k=0}^{n} a_{k,n-k}
\end{align*}
and similarly
\begin{align*}
  \sum_{k=0}^{+\infty} \sum_{n=0}^{+\infty} \abs{b_{k,n}}
  &= \sum_{k=0}^{+\infty} \sum_{j=0}^{+\infty} \abs{a_{k,j}}
  \\
  \sum_{n=0}^{+\infty} \sum_{k=0}^{+\infty} \abs{b_{k,n}}
  &= \sum_{n=0}^{+\infty} \sum_{k=0}^{n} \abs{a_{k,n-k}}.  
\end{align*}
Thus, theorem~\ref{th:scambio_somma} can be applied to exchange the sums of the
series with terms $b_{k,n}$ to obtain the result:
\[
  \sum_{k=0}^{+\infty} \sum_{j=0}^{+\infty} a_{k,j}
  = \sum_{k=0}^{+\infty}\sum_{n=0}^{+\infty} b_{k,n} 
  = \sum_{n=0}^{+\infty} \sum_{k=0}^{+\infty} b_{k,n}
  = \sum_{n=0}^{+\infty} \sum_{k=0}^{n} a_{k,n-k}.
\]
%from which~\eqref{eq:somma_Cauchy}.
\end{proof}